\documentclass[a4paper,12pt]{article}

% Pacchetti
\usepackage[utf8]{inputenc}
\usepackage{xcolor}
\usepackage{geometry}
\usepackage{listings}
\usepackage{fancyhdr}
\usepackage{titlesec}
\geometry{margin=1in}

% Configurazione codice
\lstset{
    language=C,
    basicstyle=\ttfamily\footnotesize,
    keywordstyle=\color{blue}\bfseries,
    stringstyle=\color{red},
    commentstyle=\color{gray},
    numbers=left,
    numberstyle=\tiny\color{gray},
    stepnumber=1,
    numbersep=5pt,
    tabsize=4,
    breaklines=true,
    showstringspaces=false,
    frame=single,
    captionpos=b
}

% Intestazione e piè di pagina
\pagestyle{fancy}
\fancyhead[L]{Programmazione in C}
\fancyhead[R]{\thepage}
\fancyfoot[L]{Esempi pratici}
\fancyfoot[R]{\today}

% Titoli personalizzati
\titleformat{\section}{\Large\bfseries}{\thesection}{1em}{}
\titleformat{\subsection}{\large\bfseries}{\thesubsection}{1em}{}

% Titolo
\title{\textbf{Esempi di Programmazione in C}}
\author{Studente di Informatica}
\date{\today}

\begin{document}

\maketitle

\tableofcontents
\newpage

\section{Introduzione}
Questo documento presenta esempi pratici di programmazione in C, focalizzandosi su tre argomenti principali:
\begin{itemize}
    \item Funzioni avanzate della libreria \texttt{string.h}.
    \item Operazioni avanzate sui file.
    \item Generazione di numeri casuali.
\end{itemize}

Ogni sezione contiene esempi di codice dettagliati e spiegazioni, utili per studenti e sviluppatori.

\section{Funzioni della Libreria \texttt{string.h}}
La libreria \texttt{string.h} include molte funzioni utili per la manipolazione delle stringhe.

\subsection{Funzioni Base}
Ecco alcune delle funzioni più utilizzate:
\begin{itemize}
    \item \texttt{strcat}: Concatena due stringhe.
    \item \texttt{strcmp}: Confronta due stringhe.
    \item \texttt{strlen}: Calcola la lunghezza di una stringa.
    \item \texttt{strcpy}: Copia una stringa in un'altra.
\end{itemize}

\subsection{Funzioni Avanzate}
Esempio di utilizzo di funzioni come \texttt{strstr}, \texttt{strchr} e \texttt{strcasecmp}:
\begin{lstlisting}[caption={Esempio di utilizzo avanzato di \texttt{string.h}}]
#include <stdio.h>
#include <string.h>

int main() {
    char testo[] = "C programming is fun!";
    char parola[] = "programming";

    // Cercare una sottostringa
    char *sottostringa = strstr(testo, parola);
    if (sottostringa) {
        printf("Trovata la sottostringa: %s\n", sottostringa);
    }

    // Cercare un carattere
    char *carattere = strchr(testo, 'f');
    if (carattere) {
        printf("Trovato il carattere 'f' alla posizione: %ld\n", carattere - testo);
    }

    // Confrontare stringhe ignorando maiuscole/minuscole
    if (strcasecmp("HELLO", "hello") == 0) {
        printf("Le stringhe sono uguali (case insensitive).\n");
    }

    return 0;
}
\end{lstlisting}

\section{Operazioni Avanzate sui File}
In C, i file vengono gestiti tramite la libreria \texttt{stdio.h}. Ecco alcune operazioni avanzate:

\subsection{Appendere Dati a un File}
\begin{lstlisting}[caption={Appendere dati a un file}]
#include <stdio.h>

int main() {
    FILE *fp = fopen("output.txt", "a");
    if (fp == NULL) {
        perror("Errore apertura file");
        return 1;
    }

    fprintf(fp, "Aggiungo una nuova riga al file.\n");
    fclose(fp);
    printf("Dati aggiunti al file.\n");

    return 0;
}
\end{lstlisting}

\subsection{Contare le Linee in un File}
\begin{lstlisting}[caption={Contare le linee in un file}]
#include <stdio.h>

int main() {
    FILE *fp = fopen("output.txt", "r");
    if (fp == NULL) {
        perror("Errore apertura file");
        return 1;
    }

    int linee = 0;
    char buffer[100];
    while (fgets(buffer, 100, fp) != NULL) {
        linee++;
    }
    fclose(fp);

    printf("Il file contiene %d linee.\n", linee);
    return 0;
}
\end{lstlisting}

\section{Generazione di Numeri Casuali}
La generazione di numeri casuali in C è gestita tramite la libreria \texttt{stdlib.h}.

\subsection{Esempio di Numeri Casuali}
\begin{lstlisting}[caption={Generazione di numeri casuali tra 1 e 100}]
#include <stdio.h>
#include <stdlib.h>
#include <time.h>

int main() {
    srand(time(NULL)); // Inizializza il generatore

    // Genera 5 numeri casuali tra 1 e 100
    for (int i = 0; i < 5; i++) {
        int numero = (rand() % 100) + 1;
        printf("Numero casuale: %d\n", numero);
    }

    return 0;
}
\end{lstlisting}

\section{Esempio Combinato}
Un esempio complesso che combina stringhe, file e numeri casuali:
\begin{lstlisting}[caption={Esempio combinato con stringhe, file e random}]
#include <stdio.h>
#include <stdlib.h>
#include <string.h>
#include <time.h>

int main() {
    FILE *fp = fopen("random_data.txt", "w");
    if (fp == NULL) {
        perror("Errore apertura file");
        return 1;
    }

    srand(time(NULL));
    for (int i = 0; i < 5; i++) {
        char nome[20];
        sprintf(nome, "Nome%d", i + 1);

        int numero = (rand() % 100) + 1;

        fprintf(fp, "%s - Numero: %d\n", nome, numero);
    }
    fclose(fp);

    printf("Dati scritti nel file.\n");
    return 0;
}
\end{lstlisting}

\end{document}
